%% disertasi.tex
%% Master file untuk Disertasi Doktor ITB
%%
%% Build:
%%   make watch        # auto-recompile saat menulis
%%   make build        # compile sekali
%%   make pre-submit   # cek lengkap sebelum kirim ke promotor
%%
%% Engine: LuaLaTeX (preferred) atau pdfLaTeX
%% Backend bibliografi: biber

\documentclass[final,lualatex]{itbdisertasi}

% =====================================================================
% Metadata Disertasi
% =====================================================================
\title{Prediksi Sisa Umur Pakai Bantalan Gelinding Menggunakan Arsitektur
       Hibrida Mamba-xLSTM dengan Interpretabilitas Sparse Autoencoder}
\author{Nama Lengkap Mahasiswa}
\nim{XXXXXXXXX}
\prodi{Teknik Mesin}
\kodeprodi{36101}
\fakultas{Fakultas Teknik Mesin dan Dirgantara}
\bulan{Mei}
\tahun{2026}
\promotor{Prof. Dr. Ir. Nama Promotor, M.Sc.}
\nipPromotor{19XX0000XXXXX0X0XXX}
\kopromotor{Dr. Ir. Nama Ko-Promotor, M.T.}
\nipKopromotor{19XX0000XXXXX0X0XXX}

% PDF metadata
\hypersetup{
  pdftitle={Prediksi Sisa Umur Pakai Bantalan Gelinding Menggunakan Mamba-xLSTM dan SAE},
  pdfauthor={Nama Lengkap Mahasiswa},
  pdfsubject={Disertasi Doktor ITB - Teknik Mesin},
  pdfkeywords={RUL, bantalan gelinding, Mamba, xLSTM, sparse autoencoder, interpretabilitas, pemeliharaan prediktif}
}

% Bibliografi
\addbibresource{references.bib}

% Direktori gambar
\graphicspath{{figures/}{figures/bab1/}{figures/bab2/}{figures/bab3/}{figures/bab4/}{figures/bab5/}{figures/bab6/}}

% =====================================================================
\begin{document}

% ---------------------------------------------------------------------
% Bagian Persiapan (Front Matter)
% ---------------------------------------------------------------------
\sampul

\itbfrontmatter

\halamanpengesahan
\pedomanpenggunaan

%% chapters/00-abstrak-id.tex
%% Abstrak Bahasa Indonesia — 500-800 kata, max 7 kata kunci
%% TIDAK boleh: sitasi, gambar, tabel, persamaan, footnote
%% Aturan: lihat .cursor/rules/12-abstrak-front-matter.mdc

\begin{abstrak}

Prediksi sisa umur pakai (\emph{remaining useful life}, RUL) bantalan
gelinding berperan sentral pada strategi pemeliharaan prediktif modern.
Metode mutakhir berbasis \emph{Long Short-Term Memory} (LSTM) dan
\emph{Transformer} masih menghadapi \emph{trade-off} antara akurasi pada
sekuens panjang, biaya komputasi, dan keterjelasan keputusan model
terhadap fenomena fisis degradasi bantalan. Disertasi ini mengusulkan
\emph{Mamba-xLSTM-Net}, arsitektur hibrida yang mengintegrasikan
\emph{state-space model} selektif (Mamba) untuk pemodelan dinamika
jangka panjang dengan modul \emph{xLSTM} bermemori matriks untuk
mempertahankan presisi rekuren. Arsitektur tersebut dilengkapi modul
interpretabilitas \emph{post-hoc} berbasis \emph{Top-k Sparse
Autoencoder} (SAE) yang dapat dipetakan secara empiris ke frekuensi
karakteristik bantalan (\emph{Ball Pass Frequency Outer race},
\emph{Ball Pass Frequency Inner race}, \emph{Ball Spin Frequency},
\emph{Fundamental Train Frequency}).

Eksperimen dilakukan pada dua dataset publik yaitu PHM2012 dan XJTU-SY
dengan baseline mencakup LSTM, Transformer, Informer, Mamba murni, dan
xLSTM murni. Setiap konfigurasi dilatih dengan lima inisialisasi acak
yang berbeda dan dievaluasi menggunakan metrik \emph{Root Mean Square
Error} (RMSE), \emph{Mean Absolute Error} (MAE), serta \emph{scoring
function} PHM2012. Hasil utama menunjukkan penurunan RMSE rerata
sebesar XX~\% dan MAE sebesar XX~\% dibanding baseline terkuat, dengan
keunggulan paling konsisten pada kondisi operasi beban tinggi. Studi
ablasi mengonfirmasi bahwa kombinasi Mamba dan xLSTM memberikan
kontribusi aditif: penghapusan salah satu modul menyebabkan degradasi
kinerja yang signifikan secara statistik. Analisis interpretabilitas
mengungkap bahwa fitur-fitur laten SAE yang paling sering aktif pada
fase \emph{late-life} berkorespondensi kuat dengan komponen spektral di
sekitar BPFO, BPFI, dan BSF teoretis, memberikan bukti empiris bahwa
model belajar representasi yang konsisten dengan teori klasik diagnosis
getaran bantalan.

Kontribusi utama disertasi ini mencakup: (1) arsitektur hibrida
Mamba-xLSTM-Net yang menggabungkan keunggulan \emph{state-space model}
selektif dan \emph{matrix-memory recurrent}; (2) metodologi
interpretabilitas SAE yang menghubungkan fitur laten dengan konsep
fisis frekuensi karakteristik bantalan; serta (3) bukti empiris
peningkatan kinerja prediksi RUL dengan margin signifikan pada dua
dataset publik yang digunakan secara luas oleh komunitas prognostik.
Hasil disertasi ini membuka jalan bagi pengembangan sistem prognostik
berbasis \emph{deep learning} yang lebih akurat sekaligus dapat
dipercaya oleh para insinyur pemeliharaan di industri.

\vspace{0.5cm}
\noindent
\textbf{Kata kunci:} prediksi sisa umur pakai, bantalan gelinding,
Mamba, xLSTM, sparse autoencoder, interpretabilitas, pemeliharaan
prediktif

\end{abstrak}

%% chapters/00-abstract-en.tex
%% Abstract Bahasa Inggris — content equivalen dengan abstrak ID
%% Aturan: lihat .cursor/rules/12-abstrak-front-matter.mdc

\begin{abstract*}

Remaining Useful Life (RUL) prediction for rolling element bearings
underpins modern predictive maintenance strategies. State-of-the-art
methods based on Long Short-Term Memory (LSTM) networks and
Transformers face a trade-off between long-sequence accuracy,
computational cost, and the interpretability of model decisions with
respect to the underlying physical degradation phenomena of bearings.
This dissertation proposes \emph{Mamba-xLSTM-Net}, a hybrid
architecture that integrates a selective state-space model (Mamba) for
long-range dynamics with an xLSTM module equipped with matrix memory
for precise recurrent state tracking. The architecture is augmented by
a post-hoc interpretability module based on a Top-$k$ Sparse
Autoencoder (SAE) whose latent features are mapped empirically to the
characteristic bearing frequencies (Ball Pass Frequency Outer race,
Ball Pass Frequency Inner race, Ball Spin Frequency, and Fundamental
Train Frequency).

Experiments were conducted on two public datasets, PHM2012 and
XJTU-SY, with baselines including LSTM, Transformer, Informer, vanilla
Mamba, and vanilla xLSTM. Each configuration was trained with five
random seeds and evaluated using Root Mean Square Error (RMSE), Mean
Absolute Error (MAE), and the PHM2012 scoring function. The proposed
method achieved an average RMSE reduction of XX\% and an MAE reduction
of XX\% relative to the strongest baseline, with the largest gains
under high-load operating conditions. Ablation studies confirmed that
the Mamba and xLSTM components contribute additively: removing either
module led to a statistically significant performance drop. The
interpretability analysis revealed that the SAE latent features most
frequently activated during late-life regimes correspond strongly with
spectral components near the theoretical BPFO, BPFI, and BSF,
providing empirical evidence that the model learns representations
consistent with classical bearing vibration diagnostics.

The main contributions of this dissertation are: (i) the
Mamba-xLSTM-Net hybrid architecture that combines the strengths of
selective state-space models and matrix-memory recurrent networks;
(ii) an interpretability methodology that links SAE latent features to
the physical concept of bearing characteristic frequencies; and (iii)
empirical evidence of significant performance improvement on two
widely used benchmarks. The results pave the way for deep
learning-based prognostic systems that are simultaneously more
accurate and more trustworthy for maintenance engineers in industry.

\vspace{0.5cm}
\noindent
\textbf{Keywords:} remaining useful life, rolling bearing, Mamba,
xLSTM, sparse autoencoder, interpretability, predictive maintenance

\end{abstract*}

%% chapters/00-kata-pengantar.tex
%% Kata Pengantar / Ucapan Terima Kasih
%% Panjang: 1-2 halaman, profesional, hindari informalitas berlebihan

\chapter*{KATA PENGANTAR}
\addcontentsline{toc}{chapter}{KATA PENGANTAR}

Puji syukur dipanjatkan ke hadirat Tuhan Yang Maha Esa atas segala
rahmat dan karunia-Nya sehingga penulisan disertasi yang berjudul
``\textbf{\@title}'' dapat diselesaikan sebagai salah satu syarat
memperoleh gelar Doktor pada Program Studi Doktor \@prodi, Institut
Teknologi Bandung.

Penyelesaian disertasi ini tidak lepas dari bantuan, bimbingan, dan
dukungan banyak pihak. Pada kesempatan ini, penulis menyampaikan
penghargaan dan ucapan terima kasih yang sebesar-besarnya kepada:

\begin{enumerate}
  \item \@promotor{} selaku Promotor, atas bimbingan, arahan ilmiah,
        dan kesabaran beliau selama penyelesaian disertasi ini;
  \item \@kopromotor{} selaku Ko-Promotor, atas diskusi mendalam dan
        masukan berharga pada metodologi dan analisis;
  \item Tim Penguji Disertasi yang telah memberikan saran konstruktif
        untuk perbaikan;
  \item Ketua dan staf Program Studi Doktor \@prodi{} ITB yang telah
        memfasilitasi seluruh proses akademik;
  \item Pemberi beasiswa / sponsor penelitian, atas dukungan finansial
        yang memungkinkan penelitian ini berjalan;
  \item Rekan-rekan peneliti di laboratorium, atas diskusi ilmiah dan
        kebersamaan;
  \item Keluarga tercinta, atas doa dan dukungan tanpa syarat sepanjang
        masa studi.
\end{enumerate}

Penulis menyadari bahwa disertasi ini masih jauh dari sempurna. Kritik
dan saran yang membangun sangat diharapkan untuk pengembangan lebih
lanjut. Semoga karya ini dapat memberikan kontribusi yang berarti bagi
perkembangan ilmu pengetahuan, khususnya di bidang prognostik dan
pemeliharaan prediktif berbasis \emph{deep learning}.

\vspace{1cm}
\begin{flushright}
Bandung, \@bulan{} \@tahun\\[0.5cm]
Penulis
\end{flushright}

\cleardoublepage


\renewcommand{\contentsname}{DAFTAR ISI}
\tableofcontents
\cleardoublepage

\renewcommand{\listfigurename}{DAFTAR GAMBAR}
\listoffigures
\cleardoublepage

\renewcommand{\listtablename}{DAFTAR TABEL}
\listoftables
\cleardoublepage

%% chapters/00-daftar-singkatan.tex
%% Daftar Singkatan dan Lambang
%% Wajib alfabetis. Untuk lambang matematis, kelompokkan terpisah.

\chapter*{DAFTAR SINGKATAN DAN LAMBANG}
\addcontentsline{toc}{chapter}{DAFTAR SINGKATAN DAN LAMBANG}

\section*{Singkatan}
\begin{singlespace}
\begin{tabular}{p{3cm}p{11cm}}
\textbf{Singkatan} & \textbf{Kepanjangan} \\[2pt]
\hline\\[-6pt]
BPFI  & \emph{Ball Pass Frequency Inner race} \\
BPFO  & \emph{Ball Pass Frequency Outer race} \\
BSF   & \emph{Ball Spin Frequency} \\
CNN   & \emph{Convolutional Neural Network} \\
EOL   & \emph{End of Life} \\
FPT   & \emph{First Prediction Time} \\
FTF   & \emph{Fundamental Train Frequency} \\
GRU   & \emph{Gated Recurrent Unit} \\
HI    & \emph{Health Indicator} \\
IG    & \emph{Integrated Gradients} \\
LSTM  & \emph{Long Short-Term Memory} \\
MAE   & \emph{Mean Absolute Error} \\
MAPE  & \emph{Mean Absolute Percentage Error} \\
mLSTM & \emph{matrix LSTM} (komponen xLSTM) \\
MLP   & \emph{Multi-Layer Perceptron} \\
PdM   & \emph{Predictive Maintenance} \\
PHM   & \emph{Prognostics and Health Management} \\
RMSE  & \emph{Root Mean Square Error} \\
RNN   & \emph{Recurrent Neural Network} \\
RUL   & \emph{Remaining Useful Life} (Sisa Umur Pakai) \\
SAE   & \emph{Sparse Autoencoder} \\
SHAP  & \emph{SHapley Additive exPlanations} \\
sLSTM & \emph{scalar LSTM} (komponen xLSTM) \\
SOTA  & \emph{State-of-the-Art} \\
SSM   & \emph{State-Space Model} \\
xLSTM & \emph{Extended LSTM} \\
\end{tabular}
\end{singlespace}

\vspace{1cm}
\section*{Lambang}
\begin{singlespace}
\begin{tabular}{p{3cm}p{11cm}}
\textbf{Lambang} & \textbf{Arti} \\[2pt]
\hline\\[-6pt]
$\mathbf{x}_t$ & Vektor sinyal getaran pada waktu $t$ \\
$\mathbf{h}_t$ & Vektor laten / \emph{hidden state} pada waktu $t$ \\
$\mathbf{C}_t$ & Matriks memori mLSTM pada waktu $t$ \\
$\mathbf{A}, \mathbf{B}, \mathbf{C}$ & Parameter SSM kontinu \\
$\bar{\mathbf{A}}, \bar{\mathbf{B}}$ & Parameter SSM terdiskretisasi \\
$\Delta$       & \emph{Step size} diskretisasi SSM \\
$y_i$          & Nilai RUL aktual untuk sampel ke-$i$ \\
$\hat{y}_i$    & Nilai RUL prediksi untuk sampel ke-$i$ \\
$f_r$          & Frekuensi rotasi poros (Hz) \\
$n$            & Jumlah \emph{rolling element} pada bantalan \\
$d$            & Diameter \emph{rolling element} (mm) \\
$D$            & \emph{Pitch diameter} bantalan (mm) \\
$\theta$       & Sudut kontak bantalan (rad) \\
$\lambda$      & Bobot regularisasi sparsity SAE \\
$k$            & Jumlah aktivasi yang dipertahankan pada Top-$k$ SAE \\
$\sigma(\cdot)$ & Fungsi aktivasi sigmoid \\
$\odot$        & Perkalian elemen demi elemen (Hadamard) \\
$\mathcal{L}$  & Fungsi kerugian (\emph{loss}) \\
\end{tabular}
\end{singlespace}

\cleardoublepage


% ---------------------------------------------------------------------
% Bagian Utama (Main Matter)
% ---------------------------------------------------------------------
\itbmainmatter

%% chapters/01-pendahuluan.tex
%% Bab I — Pendahuluan
%% Aturan: lihat .cursor/rules/06-bab1-pendahuluan.mdc

\chapter{Pendahuluan}
\label{bab:pendahuluan}

\section{Latar Belakang}
\label{subsec:bab1_latar}

Bantalan gelinding (\emph{rolling element bearings}) adalah komponen
mekanis yang sangat luas digunakan dalam mesin rotari modern, mulai
dari motor listrik berkecepatan tinggi hingga turbin angin
berkapasitas \si{\mega\watt}. Kegagalan bantalan menjadi penyebab
sekitar XX\,\% dari seluruh kerusakan mesin rotari di industri
\citetitb{Reference1}, sehingga kemampuan memprediksi sisa umur pakai
(\emph{remaining useful life}, RUL) bantalan secara akurat menjadi
fondasi strategi pemeliharaan prediktif (\emph{predictive maintenance}).

% [paragraph 2: tren PHM dan deep learning]
% [paragraph 3: keterbatasan SOTA — LSTM saturasi pada sekuens panjang;
%  Transformer mahal komputasi; minim interpretabilitas]
% [paragraph 4: peluang Mamba dan xLSTM, motivasi hibrida]
% [paragraph 5: pentingnya interpretabilitas — sparse autoencoder]
% [paragraph 6: posisi dan ringkasan kontribusi disertasi ini]

\section{Masalah Penelitian}
\label{subsec:bab1_masalah}

Berdasarkan latar belakang di atas, dirumuskan tiga pertanyaan
penelitian (\emph{research questions}, RQ) sebagai berikut:

\begin{description}
  \item[RQ1.] Bagaimana mendesain arsitektur \emph{deep learning} yang
        memodelkan dinamika sinyal getaran jangka panjang sekaligus
        mempertahankan presisi rekuren untuk prediksi RUL bantalan
        yang lebih akurat dibanding pendekatan SOTA berbasis
        Transformer maupun LSTM klasik?
  \item[RQ2.] Bagaimana modul interpretabilitas \emph{post-hoc} dapat
        dirancang sehingga representasi laten model RUL dapat
        dipetakan secara empiris ke konsep fisis yang dapat
        diverifikasi dengan teori klasik diagnosis getaran bantalan?
  \item[RQ3.] Sejauh mana arsitektur dan metodologi yang diusulkan
        dapat digeneralisasi lintas dataset dengan kondisi operasi
        yang berbeda?
\end{description}

\section{Tujuan Penelitian}
\label{subsec:bab1_tujuan}

Penelitian ini bertujuan untuk:

\begin{enumerate}
  \item Mengembangkan arsitektur hibrida \emph{Mamba-xLSTM-Net} yang
        menggabungkan keunggulan \emph{state-space model} selektif
        dan \emph{matrix-memory recurrent} untuk prediksi RUL
        bantalan gelinding.
  \item Merancang dan memvalidasi modul interpretabilitas berbasis
        \emph{Top-k Sparse Autoencoder} yang dapat dipetakan ke
        frekuensi karakteristik bantalan.
  \item Memvalidasi metodologi yang diusulkan secara empiris pada
        dataset publik PHM2012 dan XJTU-SY dengan baseline
        komprehensif.
\end{enumerate}

\section{Lingkup Penelitian}
\label{subsec:bab1_lingkup}

Penelitian ini dibatasi pada:

\begin{itemize}
  \item Bantalan gelinding jenis \emph{deep groove ball bearing} pada
        rezim degradasi alami (\emph{natural run-to-failure}).
  \item Sinyal getaran satu sumbu (vertikal) dengan frekuensi sampling
        \SI{25.6}{\kilo\hertz}.
  \item Dataset publik PHM2012 (FEMTO-PRONOSTIA) dan XJTU-SY.
  \item Prediksi RUL skalar (regresi tunggal); tidak mencakup
        klasifikasi mode kegagalan.
  \item Interpretabilitas pada level fitur laten dan korespondensi
        spektral; tidak mencakup interpretabilitas kausal melalui
        intervensi.
\end{itemize}

\section{Asumsi Penelitian}
\label{subsec:bab1_asumsi}

Asumsi yang digunakan:

\begin{enumerate}
  \item Degradasi bantalan diasumsikan monoton secara umum, dengan
        toleransi terhadap fluktuasi lokal.
  \item Label RUL dihitung secara linier menurun dari titik awal
        prediksi (\emph{first prediction time}, FPT) hingga \emph{end
        of life} (EOL).
  \item Karakteristik bantalan dan kondisi operasi pada setiap
        dataset dianggap homogen sesuai dokumentasi resmi dataset.
\end{enumerate}

\section{Hipotesis Penelitian}
\label{subsec:bab1_hipotesis}

\begin{description}
  \item[H1.] Kombinasi \emph{state-space model} selektif (Mamba) dan
        \emph{matrix-memory recurrent} (mLSTM) memberikan kinerja
        prediksi RUL yang lebih baik dibanding masing-masing komponen
        secara terpisah.
  \item[H2.] Modul \emph{Top-k Sparse Autoencoder} yang diterapkan
        secara \emph{post-hoc} dapat menghasilkan fitur laten yang
        konsisten dengan frekuensi karakteristik bantalan tanpa
        mendegradasi akurasi prediksi RUL.
  \item[H3.] Arsitektur yang diusulkan menunjukkan generalisasi yang
        lebih baik lintas dataset dibanding baseline berbasis
        Transformer.
\end{description}

\section{Pendekatan Penelitian}
\label{subsec:bab1_pendekatan}

Penelitian ini menggunakan pendekatan eksperimental kuantitatif
dengan urutan: studi literatur, perumusan arsitektur, implementasi
prototipe, eksperimen sistematis dengan ablasi, analisis
interpretabilitas, dan validasi lintas dataset. Detail metodologi
dijabarkan pada Bab~\ref{bab:metodologi}.

\section{Kebaruan Penelitian}
\label{subsec:bab1_kebaruan}

Disertasi ini menyumbangkan tiga kebaruan utama (kategorisasi
sesuai Pedoman ITB §V.1):

\begin{enumerate}
  \item \textbf{(Teknologi-Metodologi)} Arsitektur hibrida
        \emph{Mamba-xLSTM-Net} yang mengintegrasikan
        \emph{state-space model} selektif dengan
        \emph{matrix-memory xLSTM} untuk prediksi RUL bantalan,
        dengan mekanisme fusi yang dirancang berdasarkan sifat
        non-stasioner sinyal getaran.
  \item \textbf{(Teknologi-Metodologi)} Skema interpretabilitas
        \emph{Top-k Sparse Autoencoder} yang bekerja secara
        \emph{post-hoc} terhadap representasi laten model RUL,
        dilengkapi prosedur kuantitatif untuk memetakan fitur SAE
        ke frekuensi karakteristik fisik bantalan.
  \item \textbf{(Output)} Bukti empiris bahwa model \emph{deep
        learning} yang dilatih untuk RUL mempelajari representasi
        laten yang berkorespondensi dengan teori klasik diagnosis
        getaran bantalan, disertai peningkatan kinerja SOTA pada
        dua dataset publik.
\end{enumerate}

\section{Sistematika Penulisan}
\label{subsec:bab1_sistematika}

Disertasi ini disusun ke dalam enam bab. \textbf{Bab~I} memuat latar
belakang, masalah, tujuan, lingkup, asumsi, hipotesis, pendekatan,
kebaruan, dan sistematika penulisan. \textbf{Bab~II} menyajikan
tinjauan pustaka secara kronologis terkait metode prediksi RUL
bantalan dan perkembangan arsitektur \emph{deep learning}.
\textbf{Bab~III} memuat dasar teori meliputi formulasi RUL, LSTM,
xLSTM (sLSTM dan mLSTM), \emph{state-space model} dan Mamba,
\emph{Sparse Autoencoder}, serta metode atribusi.
\textbf{Bab~IV} menjabarkan metodologi penelitian, termasuk
arsitektur yang diusulkan, dataset, protokol pelatihan, dan
rancangan eksperimen. \textbf{Bab~V} menyajikan hasil eksperimen,
studi ablasi, dan analisis interpretabilitas. \textbf{Bab~VI}
menutup disertasi dengan kesimpulan dan saran pekerjaan lanjut.

%% chapters/02-tinjauan-pustaka.tex
%% Bab II — Tinjauan Pustaka
%% Catatan: Tinjauan Pustaka != Dasar Teori. Bab ini menyajikan EVOLUSI
%% pekerjaan terdahulu secara kronologis dan menempatkan disertasi ini
%% terhadap state-of-the-art. Aturan: lihat .cursor/rules/07-bab2-tinjauan-pustaka.mdc

\chapter{Tinjauan Pustaka}
\label{bab:tinjauan-pustaka}

Bab ini menyajikan tinjauan kritis atas pekerjaan terdahulu yang
relevan dengan disertasi ini. Pembahasan disusun secara kronologis
dari metode prediksi RUL klasik berbasis fisika hingga pendekatan
\emph{deep learning} mutakhir, serta menempatkan kontribusi disertasi
ini terhadap \emph{state-of-the-art}.

\section{Domain Prognostik Bantalan Gelinding}
\label{subsec:bab2_domain}

% Konteks sejarah singkat: ISO standards, condition monitoring, evolusi
% dari corrective ke predictive maintenance. Frekuensi karakteristik
% (BPFO/BPFI/BSF/FTF) sebagai fondasi diagnostik. Fitur statistik
% klasik (RMS, kurtosis) dalam pemantauan kondisi.
\dots

\section{Evolusi Metode Prediksi RUL}
\label{subsec:bab2_evolusi}

\subsection{Metode Berbasis Model Fisika}
\label{subsec:bab2_fisika}

% Paris law, Forman law, model akumulasi kerusakan. Kelebihan: dapat
% diinterpretasi. Keterbatasan: butuh model fisika spesifik bantalan
% dan kondisi operasi yang akurat.
\dots

\subsection{Metode Berbasis Data — Pra Era \emph{Deep Learning}}
\label{subsec:bab2_data_pre_dl}

% Hidden Markov Model, Particle Filter, SVR, Kalman filter, ARMA, dan
% kombinasi feature engineering manual + classifier dangkal.
\dots

\subsection{Era \emph{Deep Learning}: RNN, LSTM, dan GRU}
\label{subsec:bab2_dl_lstm}

% Pekerjaan-pekerjaan kunci 2016--2020 yang menerapkan LSTM/GRU pada
% PHM2012 dan dataset bearing lainnya. Pencapaian dan keterbatasan.
\dots

\subsection{Era Transformer dan Variannya}
\label{subsec:bab2_transformer}

% Self-attention untuk time series, Informer, Autoformer, FEDformer,
% PatchTST. Trade-off akurasi vs biaya komputasi.
\dots

\section{State-of-the-Art Terkini: xLSTM dan Mamba}
\label{subsec:bab2_sota}

\subsection{xLSTM dan \emph{Matrix Memory}}
\label{subsec:bab2_xlstm}

% Pekerjaan Beck dkk. (2024) yang memperkenalkan xLSTM dengan komponen
% sLSTM dan mLSTM. Aplikasi xLSTM pada time series.
\dots

\subsection{Mamba dan \emph{Selective State-Space Models}}
\label{subsec:bab2_mamba}

% Pekerjaan Gu \& Dao (2023) tentang Mamba. Adopsi pada time series
% forecasting dan prognostik.
\dots

\section{Interpretabilitas pada Model \emph{Deep Learning}}
\label{subsec:bab2_interpret}

% SHAP, LIME, Integrated Gradients pada model time series. Sparse
% Autoencoder di NLP/vision (Anthropic features, OpenAI scaling).
% Adopsi terbatas SAE pada prognostik.
\dots

\section{Posisi Disertasi dan Celah Penelitian}
\label{subsec:bab2_posisi}

Berdasarkan tinjauan di atas, dapat diidentifikasi tiga celah utama
yang akan dialamatkan oleh disertasi ini:

\begin{enumerate}
  \item Belum ada arsitektur yang secara eksplisit menggabungkan
        \emph{state-space model} selektif dengan \emph{matrix-memory
        recurrent} untuk prediksi RUL bantalan.
  \item Aplikasi \emph{Sparse Autoencoder} pada prognostik bantalan
        masih sangat terbatas, dan belum ada studi yang memetakan
        fitur SAE ke frekuensi karakteristik fisik.
  \item Studi generalisasi lintas dataset masih jarang, sehingga
        klaim SOTA pada satu dataset belum tentu konsisten lintas
        kondisi operasi.
\end{enumerate}

% Tabel komparatif: 8-12 baris pekerjaan terkait, kolom: tahun, metode,
% dataset, metrik, kontribusi, keterbatasan. Letakkan baris terakhir
% sebagai disertasi ini.
\begin{table}[ht]
\centering
\caption{Posisi disertasi terhadap pekerjaan terdahulu yang representatif.}
\label{tab:bab2_posisi}
\small
\begin{tabular}{@{}p{2.2cm}p{2cm}p{2cm}p{2.2cm}p{4cm}@{}}
\toprule
Studi & Arsitektur & Dataset & Interpret. & Keterbatasan \\
\midrule
\dots & \dots & \dots & \dots & \dots \\
\midrule
\textbf{Disertasi ini} & Mamba-xLSTM-Net & PHM2012, XJTU-SY & SAE + SHAP + IG & --- \\
\bottomrule
\end{tabular}
\end{table}

%% chapters/03-dasar-teori.tex
%% Bab III — Dasar Teori
%% Aturan: lihat .cursor/rules/08-bab3-dasar-teori.mdc
%% Berisi formulasi matematis lengkap; Bab IV merujuk persamaan-persamaan di sini.

\chapter{Dasar Teori}
\label{bab:dasar-teori}

Bab ini memuat fondasi teoretis yang menjadi dasar metodologi pada
Bab~\ref{bab:metodologi}. Disajikan formulasi matematis dari prediksi
RUL, varian arsitektur LSTM (sLSTM, mLSTM), \emph{state-space model}
dan Mamba, \emph{attention}, \emph{Sparse Autoencoder}, metode
atribusi, serta metrik evaluasi.

\section{Formulasi Prediksi RUL}
\label{subsec:bab3_rul}

\subsection{Definisi RUL}
\label{subsec:bab3_def_rul}

Diberikan deret waktu sinyal getaran $\{\mathbf{x}_t\}_{t=1}^{T}$
dengan $\mathbf{x}_t \in \mathbb{R}^{d_{\text{in}}}$, RUL pada waktu
$t$ didefinisikan sebagai
\begin{equation}
y_t = T_{\text{EOL}} - t,
\label{eq:bab3_rul_def}
\end{equation}
dengan $T_{\text{EOL}}$ adalah waktu \emph{end-of-life} bantalan.
Dalam praktik, RUL dinormalisasi ke selang $[0,1]$.

\subsection{Skema Pelabelan}
\label{subsec:bab3_label}

% Linear, piecewise (constant healthy phase + linear degrading phase),
% exponential. Disertasi ini akan menggunakan pendekatan piecewise
% (lihat Bab IV).
\dots

\section{Long Short-Term Memory dan Variannya}
\label{subsec:bab3_lstm}

\subsection{LSTM Klasik}
\label{subsec:bab3_lstm_classic}

% Persamaan gating LSTM original (Hochreiter & Schmidhuber, 1997).
\begin{align}
\mathbf{f}_t &= \sigma(\mathbf{W}_f \mathbf{x}_t + \mathbf{U}_f \mathbf{h}_{t-1} + \mathbf{b}_f) \label{eq:bab3_lstm_f} \\
\mathbf{i}_t &= \sigma(\mathbf{W}_i \mathbf{x}_t + \mathbf{U}_i \mathbf{h}_{t-1} + \mathbf{b}_i) \\
\mathbf{o}_t &= \sigma(\mathbf{W}_o \mathbf{x}_t + \mathbf{U}_o \mathbf{h}_{t-1} + \mathbf{b}_o) \\
\tilde{\mathbf{c}}_t &= \tanh(\mathbf{W}_c \mathbf{x}_t + \mathbf{U}_c \mathbf{h}_{t-1} + \mathbf{b}_c) \\
\mathbf{c}_t &= \mathbf{f}_t \odot \mathbf{c}_{t-1} + \mathbf{i}_t \odot \tilde{\mathbf{c}}_t \\
\mathbf{h}_t &= \mathbf{o}_t \odot \tanh(\mathbf{c}_t)
\end{align}

\subsection{sLSTM (\emph{scalar} LSTM)}
\label{subsec:bab3_slstm}

% Pekerjaan Beck dkk. (2024). Stabilizer + exponential gates.
\dots

\subsection{mLSTM (\emph{matrix} LSTM)}
\label{subsec:bab3_mlstm}

% Persamaan mLSTM dengan matrix memory.
\begin{equation}
\mathbf{C}_t = \mathbf{f}_t \mathbf{C}_{t-1} + \mathbf{i}_t \mathbf{v}_t \mathbf{k}_t^{\top}
\label{eq:bab3_mlstm_C}
\end{equation}
\dots

\section{\emph{State-Space Models} dan Mamba}
\label{subsec:bab3_ssm}

\subsection{SSM Kontinu}
\label{subsec:bab3_ssm_kontinu}

\begin{align}
\dot{\mathbf{h}}(t) &= \mathbf{A} \mathbf{h}(t) + \mathbf{B} x(t) \\
y(t) &= \mathbf{C} \mathbf{h}(t)
\label{eq:bab3_ssm_kontinu}
\end{align}

\subsection{Diskretisasi Zero-Order Hold}
\label{subsec:bab3_ssm_diskret}

\begin{align}
\bar{\mathbf{A}} &= \exp(\Delta \mathbf{A}) \\
\bar{\mathbf{B}} &= (\Delta \mathbf{A})^{-1}(\exp(\Delta \mathbf{A}) - \mathbf{I}) \cdot \Delta \mathbf{B}
\label{eq:bab3_zoh}
\end{align}

\subsection{Mamba: Selektivitas dan \emph{Selective Scan}}
\label{subsec:bab3_mamba}

% Inovasi Mamba: $\mathbf{B}, \mathbf{C}, \Delta$ tergantung input,
% sehingga selektif. Selective scan algorithm.
\dots

\section{Mekanisme \emph{Attention}}
\label{subsec:bab3_attn}

\begin{equation}
\text{Attn}(\mathbf{Q},\mathbf{K},\mathbf{V}) = \mathrm{softmax}\!\left(\frac{\mathbf{Q}\mathbf{K}^{\top}}{\sqrt{d_k}}\right)\mathbf{V}
\label{eq:bab3_attn}
\end{equation}

\section{\emph{Sparse Autoencoder}}
\label{subsec:bab3_sae}

\subsection{Formulasi Dasar}
\label{subsec:bab3_sae_basic}

\begin{align}
\mathbf{z}_t &= \text{ReLU}(\mathbf{W}_e \mathbf{h}_t + \mathbf{b}_e) \\
\hat{\mathbf{h}}_t &= \mathbf{W}_d \mathbf{z}_t + \mathbf{b}_d \\
\mathcal{L}_{\text{SAE}} &= \|\mathbf{h}_t - \hat{\mathbf{h}}_t\|_2^2 + \lambda \cdot \Omega(\mathbf{z}_t)
\label{eq:bab3_sae_loss}
\end{align}

\subsection{Top-$k$ Sparsity}
\label{subsec:bab3_topk}

\begin{equation}
\mathbf{z}_t^{(k)} = \mathrm{TopK}_k(\mathbf{z}_t)
\label{eq:bab3_topk}
\end{equation}

\section{Metode Atribusi: SHAP dan \emph{Integrated Gradients}}
\label{subsec:bab3_attribution}

\subsection{SHAP}
\label{subsec:bab3_shap}
\dots

\subsection{Integrated Gradients}
\label{subsec:bab3_ig}

\begin{equation}
\mathrm{IG}_i(\mathbf{x}) = (x_i - x'_i) \int_0^1 \frac{\partial f(\mathbf{x}' + \alpha(\mathbf{x}-\mathbf{x}'))}{\partial x_i}\, d\alpha
\label{eq:bab3_ig}
\end{equation}

\section{Frekuensi Karakteristik Bantalan}
\label{subsec:bab3_freq_char}

% Persamaan BPFO, BPFI, BSF, FTF. Lihat .cursor/rules/15-domain-rul-bearings.mdc.
\dots

\section{Metrik Evaluasi}
\label{subsec:bab3_metrik}

\begin{align}
\text{RMSE} &= \sqrt{\frac{1}{N}\sum_{i=1}^{N}(y_i - \hat{y}_i)^2} \label{eq:bab3_rmse} \\
\text{MAE}  &= \frac{1}{N}\sum_{i=1}^{N}|y_i - \hat{y}_i| \label{eq:bab3_mae}
\end{align}

\section{Notasi yang Digunakan}
\label{subsec:bab3_notasi}

% Tabel notasi terpadu — wajib ada untuk konsistensi lintas bab.
\begin{table}[ht]
\centering
\caption{Notasi terpadu yang digunakan dalam disertasi.}
\label{tab:bab3_notasi}
\small
\begin{tabular}{@{}p{2.5cm}p{11cm}@{}}
\toprule
Notasi & Arti \\
\midrule
$\mathbf{x}_t$ & Vektor input pada waktu $t$, $\mathbf{x}_t \in \mathbb{R}^{d_{\text{in}}}$ \\
$\mathbf{h}_t$ & Vektor laten / hidden state, $\mathbf{h}_t \in \mathbb{R}^{d_{\text{model}}}$ \\
$\mathbf{C}_t$ & Matriks memori mLSTM, $\mathbf{C}_t \in \mathbb{R}^{d \times d}$ \\
$\Delta$       & Step size diskretisasi SSM \\
$\mathbf{z}_t$ & Aktivasi laten SAE \\
$\lambda, k$   & Parameter regularisasi dan Top-$k$ SAE \\
\bottomrule
\end{tabular}
\end{table}

%% chapters/04-metodologi.tex
%% Bab IV — Metodologi
%% Aturan: lihat .cursor/rules/09-bab4-metodologi.mdc
%% Bab ini adalah KONTRIBUSI UTAMA — dijabarkan secara teknis dan reproducible.

\chapter{Metodologi Penelitian}
\label{bab:metodologi}

Bab ini menjabarkan metodologi penelitian secara teknis dan
reproducible: arsitektur usulan \emph{Mamba-xLSTM-Net}, dataset yang
digunakan, ekstraksi fitur, modul interpretabilitas, protokol
pelatihan, rancangan eksperimen, metrik evaluasi, dan lingkungan
komputasi.

\section{Kerangka Penelitian}
\label{subsec:bab4_kerangka}

% Gambar IV.1 — diagram alur penelitian end-to-end
\begin{figure}[ht]
\centering
% \includegraphics[width=0.9\textwidth]{figures/bab4/kerangka.pdf}
\caption{Diagram alur penelitian \emph{end-to-end} dari sinyal getaran
mentah hingga prediksi RUL dan analisis interpretabilitas.}
\label{fig:bab4_kerangka}
\end{figure}

Kerangka penelitian yang diusulkan ditunjukkan pada
Gambar~\ref{fig:bab4_kerangka}. Sinyal getaran mentah dari dataset
(Subbab~\ref{subsec:bab4_dataset}) diproses melalui ekstraksi fitur
multi-domain (Subbab~\ref{subsec:bab4_hi}), kemudian dimasukkan ke
arsitektur Mamba-xLSTM-Net (Subbab~\ref{subsec:bab4_arsitektur}) untuk
menghasilkan prediksi RUL. Modul interpretabilitas
(Subbab~\ref{subsec:bab4_interp}) bekerja secara \emph{post-hoc}
terhadap representasi laten model.

\section{Dataset}
\label{subsec:bab4_dataset}

\subsection{PHM2012 / FEMTO-PRONOSTIA}
\label{subsec:bab4_phm2012}
% Spesifikasi PHM2012 dengan tabel sesuai .cursor/rules/15-domain-rul-bearings.mdc §C.1
\dots

\subsection{XJTU-SY}
\label{subsec:bab4_xjtusy}
% Spesifikasi XJTU-SY dengan tabel sesuai .cursor/rules/15-domain-rul-bearings.mdc §C.2
\dots

\subsection{Praproses dan Segmentasi}
\label{subsec:bab4_praproses}
% Window size, overlap, normalisasi (fit hanya pada train), bearing-wise split
\dots

\section{Ekstraksi Fitur dan \emph{Health Indicator}}
\label{subsec:bab4_hi}

% Tabel 16 fitur (9 time-domain + 7 frequency-domain). Persamaan HI.
\dots

\section{Arsitektur Usulan: Mamba-xLSTM-Net}
\label{subsec:bab4_arsitektur}

% Gambar arsitektur dengan dimensi tensor di tiap blok.
\begin{figure}[ht]
\centering
% \includegraphics[width=0.95\textwidth]{figures/bab4/arsitektur.pdf}
\caption{Arsitektur \emph{Mamba-xLSTM-Net} yang diusulkan. Dimensi
tensor pada tiap blok ditandai dengan notasi $B \times T \times d$.}
\label{fig:bab4_arsitektur}
\end{figure}

\subsection{Modul Encoder Mamba}
\label{subsec:bab4_mamba}
\dots

\subsection{Modul Recurrent xLSTM}
\label{subsec:bab4_xlstm}
\dots

\subsection{Mekanisme Fusi}
\label{subsec:bab4_fusi}
\dots

\subsection{Kepala Regresi RUL}
\label{subsec:bab4_head}
\dots

\section{Modul Interpretabilitas}
\label{subsec:bab4_interp}

\subsection{Top-$k$ Sparse Autoencoder}
\label{subsec:bab4_sae}
\dots

\subsection{Atribusi Fitur: SHAP dan Integrated Gradients}
\label{subsec:bab4_shap_ig}
\dots

\section{Protokol Pelatihan}
\label{subsec:bab4_protokol}

\subsection{Fungsi Kerugian}
\label{subsec:bab4_loss}
\dots

\subsection{Optimizer dan Penjadwal Laju Belajar}
\label{subsec:bab4_optim}
\dots

\subsection{Regularisasi dan \emph{Early Stopping}}
\label{subsec:bab4_regul}
\dots

\section{Rancangan Eksperimen dan Ablasi}
\label{subsec:bab4_eksperimen}

% Daftar terstruktur eksperimen E1, E2, ... lihat aturan
% .cursor/rules/09-bab4-metodologi.mdc §H

\subsection*{Eksperimen E1 — Perbandingan SOTA}
\textbf{Tujuan:} Menjawab RQ1.\\
\textbf{Setup:} \dots\\
\textbf{Output:} Tabel V.1, Gambar V.1.

\subsection*{Eksperimen E2 — Ablasi Komponen Arsitektur}
\textbf{Tujuan:} Memvalidasi kontribusi tiap modul.\\
\textbf{Setup:} \dots

\subsection*{Eksperimen E3 — Ablasi Hyperparameter}
\dots

\subsection*{Eksperimen E4 — Analisis Interpretabilitas}
\dots

\subsection*{Eksperimen E5 — Generalisasi Lintas Dataset}
\dots

\section{Metrik Evaluasi}
\label{subsec:bab4_metrik}
% Persamaan RMSE, MAE, dan PHM2012 scoring function (asimetris)
\dots

\section{Lingkungan Komputasi}
\label{subsec:bab4_komputasi}
% Hardware, OS, versi PyTorch / mamba-ssm / xlstm, durasi training rata-rata
\dots

%% chapters/05-hasil-pembahasan.tex
%% Bab V — Hasil dan Pembahasan
%% Aturan: lihat .cursor/rules/10-bab5-hasil-pembahasan.mdc

\chapter{Hasil dan Pembahasan}
\label{bab:hasil}

Bab ini menyajikan hasil eksperimen yang dirancang pada
Subbab~\ref{subsec:bab4_eksperimen}. Setiap eksperimen dilaporkan
dalam sub-subbab tersendiri yang dilengkapi dengan tabel atau gambar
beserta diskusi naratif yang mengaitkan temuan dengan pertanyaan
penelitian dan hipotesis.

\section{Hasil Eksperimen Utama}
\label{subsec:bab5_main}

\subsection{Perbandingan dengan SOTA — XJTU-SY}
\label{subsec:bab5_main_xjtusy}

% Tabel V.1: Perbandingan kinerja Mamba-xLSTM-Net vs baseline pada XJTU-SY.
% Diskusi: pembacaan tabel, pembanding terkuat, hubungan ke H1.
\dots

\subsection{Perbandingan dengan SOTA — PHM2012}
\label{subsec:bab5_main_phm2012}
\dots

\subsection{Visualisasi Kurva Prediksi RUL}
\label{subsec:bab5_kurva}
\dots

\section{Studi Ablasi}
\label{subsec:bab5_ablasi}

\subsection{Ablasi Komponen Arsitektur}
\label{subsec:bab5_ablasi_arch}
% Tabel ablasi: -Mamba, -mLSTM, -sLSTM, -fusi, dst.
\dots

\subsection{Ablasi Hyperparameter}
\label{subsec:bab5_ablasi_hp}
% Plot tren metrik vs hyperparameter (dimensi state, jumlah head, k SAE)
\dots

\section{Analisis Interpretabilitas}
\label{subsec:bab5_interp}

\subsection{Atribusi Fitur SHAP dan Integrated Gradients}
\label{subsec:bab5_shap}
% Bar plot atribusi global + plot temporal early/mid/late life
\dots

\subsection{Sparse Autoencoder: Konsep yang Ditemukan}
\label{subsec:bab5_sae_concepts}
% Distribusi sparsity, maximum-activating examples, interpretasi 3-5 fitur SAE
\dots

\subsection{Korespondensi dengan Frekuensi Karakteristik Bantalan}
\label{subsec:bab5_korespondensi}
% Spektrum + garis vertikal BPFO/BPFI/BSF teoretis + heatmap aktivasi SAE
\dots

\section{Generalisasi Lintas Dataset}
\label{subsec:bab5_generalisasi}
% Train PHM2012 -> test XJTU-SY, dan sebaliknya
\dots

\section{Pembahasan Lebih Lanjut}
\label{subsec:bab5_diskusi}

\subsection{Kekuatan dan Kelemahan Metode Usulan}
\label{subsec:bab5_pro_kontra}
\dots

\subsection{Posisi terhadap \emph{State-of-the-Art}}
\label{subsec:bab5_posisi}
\dots

\subsection{Implikasi Praktis}
\label{subsec:bab5_implikasi}
\dots

\section{Keterbatasan Penelitian}
\label{subsec:bab5_keterbatasan}

% Daftar jujur keterbatasan yang akan menjadi seed untuk Saran di Bab VI:
% - Validasi terbatas pada bantalan gelinding deep groove ball;
% - Asumsi degradasi monoton;
% - Kondisi operasi laboratorium;
% - Klaim interpretabilitas berbasis korelasi statistik (bukan kausal);
% - Komputasi pelatihan masih relatif mahal.
\dots

%% chapters/06-kesimpulan.tex
%% Bab VI — Kesimpulan dan Saran
%% Aturan: lihat .cursor/rules/11-bab6-kesimpulan.mdc

\chapter{Kesimpulan dan Saran}
\label{bab:kesimpulan}

\section{Kesimpulan}
\label{sec:bab6_kesimpulan}

Disertasi ini telah meneliti prediksi sisa umur pakai (RUL) bantalan
gelinding menggunakan arsitektur hibrida \emph{Mamba-xLSTM-Net} dengan
modul interpretabilitas berbasis \emph{Top-k Sparse Autoencoder}. Tiga
pertanyaan penelitian yang dirumuskan pada
Subbab~\ref{subsec:bab1_masalah} dijawab sebagai berikut.

\subsection{Jawaban atas Pertanyaan Penelitian}
\label{subsec:bab6_jawaban_rq}

\paragraph{Pertanyaan penelitian pertama (RQ1)}
mempertanyakan apakah kombinasi \emph{state-space model} selektif
(Mamba) dengan \emph{matrix-memory recurrent} (mLSTM) memberikan
kinerja prediksi RUL yang lebih baik daripada masing-masing komponen
secara terpisah maupun arsitektur Transformer. Hasil eksperimen
(Subbab~\ref{subsec:bab5_main_xjtusy} dan
Subbab~\ref{subsec:bab5_main_phm2012}) menunjukkan penurunan RMSE
sebesar XX~\% pada XJTU-SY dan XX~\% pada PHM2012 dibanding baseline
terkuat. RQ1 dijawab \emph{ya}, dengan margin signifikan secara
statistik. Hasil ini mendukung Hipotesis~H1.

\paragraph{Pertanyaan penelitian kedua (RQ2)}
% Jawaban RQ2: skema interpretabilitas SAE + pemetaan ke BPFO/BPFI/BSF.
\dots

\paragraph{Pertanyaan penelitian ketiga (RQ3)}
% Jawaban RQ3: generalisasi lintas dataset.
\dots

\subsection{Kontribusi Disertasi}
\label{subsec:bab6_kontribusi}

Disertasi ini menyumbangkan empat kontribusi utama bagi bidang
prognostik bantalan gelinding:

\begin{enumerate}
  \item Mengusulkan arsitektur hibrida \emph{Mamba-xLSTM-Net} yang
        mengintegrasikan \emph{state-space model} selektif dan
        \emph{matrix-memory recurrent} dengan mekanisme fusi yang
        diturunkan dari sifat non-stasioner sinyal getaran.
  \item Memformulasikan modul interpretabilitas \emph{post-hoc}
        berbasis \emph{Top-k Sparse Autoencoder} yang bekerja
        terhadap representasi laten model RUL tanpa mendegradasi
        akurasi prediksi.
  \item Membangun korespondensi empiris antara fitur laten SAE dan
        frekuensi karakteristik bantalan (BPFO, BPFI, BSF, FTF) yang
        memberikan bukti bahwa model belajar representasi konsisten
        dengan teori klasik diagnosis getaran.
  \item Menyediakan \emph{pipeline reproducible} berikut studi ablasi
        komprehensif pada dua dataset publik (PHM2012 dan XJTU-SY)
        dengan baseline yang luas.
\end{enumerate}

\subsection{Implikasi}
\label{subsec:bab6_implikasi}

% Implikasi keilmuan: kontribusi terhadap literatur RUL dan deep learning.
% Implikasi praktis: penerapan untuk industri pemeliharaan.
% Implikasi metodologis: pelajaran tentang menggabungkan SSM dengan recurrent.
\dots

\section{Saran Pekerjaan Lanjut}
\label{sec:bab6_saran}

\subsection{Pengembangan Metodologis}
\label{subsec:bab6_saran_metode}

\paragraph{Validasi kausal terhadap konsep SAE.}
Klaim interpretabilitas dalam disertasi ini didasarkan pada korelasi
statistik antara aktivasi SAE dan komponen spektral. Pekerjaan lanjut
yang menjanjikan adalah validasi \emph{causal} melalui intervensi:
secara sengaja menambahkan komponen frekuensi karakteristik buatan ke
sinyal, lalu memeriksa apakah fitur SAE yang relevan menyala.

\paragraph{Kuantifikasi ketidakpastian prediksi.}
% Bayesian deep learning, MC Dropout, deep ensembles untuk RUL.
\dots

\subsection{Perluasan Domain Aplikasi}
\label{subsec:bab6_saran_domain}

\paragraph{Pengujian pada bantalan jenis lain.}
% Spherical roller, tapered roller, magnetic bearing.
\dots

\paragraph{Pengujian pada kondisi industri sebenarnya.}
% Data noise tinggi, beban variabel, kondisi tidak terkontrol.
\dots

\subsection{Aspek Implementasi}
\label{subsec:bab6_saran_impl}

\paragraph{Kompresi model untuk \emph{edge deployment}.}
% Pruning, quantization untuk on-device inference real-time.
\dots

\paragraph{Integrasi dengan sistem PdM industri.}
% PLC, SCADA, MQTT pipelines.
\dots


% ---------------------------------------------------------------------
% Daftar Pustaka
% ---------------------------------------------------------------------
\cleardoublepage
\printbibliography[title={DAFTAR PUSTAKA},heading=bibintoc]

% ---------------------------------------------------------------------
% Lampiran
% ---------------------------------------------------------------------
\cleardoublepage
\appendix
\renewcommand{\thechapter}{\Alph{chapter}}
\titleformat{\chapter}[display]
  {\normalfont\fontsize{14}{18}\bfseries\centering}
  {\MakeUppercase{\appendixname}\ \thechapter}
  {12pt}
  {\MakeUppercase{#1}}

%% lampiran/A-detail-fitur.tex
%% Lampiran A — Detail Fitur Time-Domain dan Frequency-Domain

\chapter{Detail Fitur Time-Domain dan Frequency-Domain}
\label{lamp:fitur}

Lampiran ini memuat formulasi lengkap dari 16 fitur yang digunakan
sebagai input pada arsitektur \emph{Mamba-xLSTM-Net}, mencakup 9
fitur \emph{time-domain} dan 7 fitur \emph{frequency-domain}.

\section{Fitur Time-Domain}
\label{lamp:fitur_time}

% Tabel 9 fitur time-domain dengan formula:
% Mean, RMS, Std, Variance, Crest factor, Impulse factor, Shape factor,
% Skewness, Kurtosis. Lihat .cursor/rules/15-domain-rul-bearings.mdc §D.1
\dots

\section{Fitur Frequency-Domain}
\label{lamp:fitur_freq}

% Tabel 7 fitur frequency-domain. Lihat
% .cursor/rules/15-domain-rul-bearings.mdc §D.2
\dots

%% lampiran/B-hyperparameter-baseline.tex
%% Lampiran B — Hyperparameter Lengkap untuk Reproducibility Baseline

\chapter{Hyperparameter Lengkap Baseline}
\label{lamp:hparam}

Lampiran ini memuat hyperparameter lengkap untuk seluruh baseline yang
dibandingkan pada Bab~\ref{bab:hasil}, demi memastikan
\emph{reproducibility}. Untuk setiap baseline, dicatat sumber
implementasi (paper / repositori GitHub / implementasi sendiri),
versi, dan deviasi dari pengaturan asli (jika ada).

\section{LSTM}
\label{lamp:hparam_lstm}
\dots

\section{Transformer}
\label{lamp:hparam_transformer}
\dots

\section{Informer}
\label{lamp:hparam_informer}
\dots

\section{Mamba (Murni)}
\label{lamp:hparam_mamba}
\dots

\section{xLSTM (Murni)}
\label{lamp:hparam_xlstm}
\dots

\section{Mamba-xLSTM-Net (Usulan)}
\label{lamp:hparam_proposed}
% Salinan Tabel IV.x untuk kemudahan rujukan
\dots

%% lampiran/C-kode-implementasi.tex
%% Lampiran C — Tautan Kode Implementasi dan Reproducibility

\chapter{Tautan Kode Implementasi dan Reproducibility}
\label{lamp:kode}

Seluruh kode implementasi, konfigurasi eksperimen, dan \emph{checkpoint}
model yang digunakan dalam disertasi ini tersedia secara publik untuk
mendukung \emph{reproducibility}.

\section{Repositori}
\label{lamp:kode_repo}

\begin{itemize}
  \item Repositori utama: \url{https://github.com/USERNAME/disertasi-mambaxlstm}
  \item Lisensi: MIT (kode), CC-BY-4.0 (dokumentasi)
  \item \emph{Commit hash} yang sesuai dengan hasil pada Bab~V:
        \texttt{XXXXXXX}
\end{itemize}

\section{Cara Menjalankan}
\label{lamp:kode_run}

% Snippet command-line untuk reproduksi eksperimen E1, E2, dst.
\begin{verbatim}
# Setup environment
conda env create -f environment.yml
conda activate disertasi

# Reproduksi eksperimen E1 (perbandingan SOTA pada XJTU-SY)
python scripts/run_experiment.py --config configs/E1_xjtusy.yaml --seed 42

# Reproduksi seluruh eksperimen (5 seed × 5 eksperimen)
bash scripts/run_all.sh
\end{verbatim}

\section{Daftar Versi Pustaka}
\label{lamp:kode_versi}

% Salin dari requirements.txt / environment.yml
\dots


\end{document}
